
%%%%%%%%%%%%%%%%%%%%%%%%%%%%%%%%%%%%%%%%%%%%%%%%%%%%%%%%%%%%%%%%%%%%%%%%%%%%%%%%%%%%%%%
%%%%%%%%%%%%%%%%%%%%%%%%%%%%%%%%%%%%%%%%%%%%%%%%%%%%%%%%%%%%%%%%%%%%%%%%%%%%%%%%%%%%%%%
% 
% This top part of the document is called the 'preamble'.  Modify it with caution!
%
% The real document starts below where it says 'The main document starts here'.

\documentclass[12pt]{article}

\usepackage{amssymb,amsmath,amsthm}
\usepackage[top=1in, bottom=1in, left=1.25in, right=1.25in]{geometry}
\usepackage{fancyhdr}
\usepackage{enumerate}
\usepackage{listings}
\usepackage{graphicx}
\usepackage{float}
\usepackage{multicol}
% Comment the following line to use TeX's default font of Computer Modern.
\usepackage{times,txfonts}
\usepackage{mwe}
\usepackage{caption}
\usepackage{subcaption}

\usepackage{tikz}
\def\checkmark{\tikz\fill[scale=0.4](0,.35) -- (.25,0) -- (1,.7) -- (.25,.15) -- cycle;} 



\makeatletter
\renewcommand*\env@matrix[1][*\c@MaxMatrixCols c]{%
  \hskip -\arraycolsep
  \let\@ifnextchar\new@ifnextchar
  \array{#1}}
\makeatother

\newtheoremstyle{homework}% name of the style to be used
  {18pt}% measure of space to leave above the theorem. E.g.: 3pt
  {12pt}% measure of space to leave below the theorem. E.g.: 3pt
  {}% name of font to use in the body of the theorem
  {}% measure of space to indent
  {\bfseries}% name of head font
  {:}% punctuation between head and body
  {2ex}% space after theorem head; " " = normal interword space
  {}% Manually specify head
\theoremstyle{homework} 

% Set up an Exercise environment and a Solution label.
\newtheorem*{exercisecore}{Exercise \@currentlabel}
\newenvironment{exercise}[1]
{\def\@currentlabel{#1}\exercisecore}
{\endexercisecore}

\newcommand{\localhead}[1]{\par\smallskip\noindent\textbf{#1}\nobreak\\}%
\newcommand\solution{\localhead{Solution:}}

%%%%%%%%%%%%%%%%%%%%%%%%%%%%%%%%%%%%%%%%%%%%%%%%%%%%%%%%%%%%%%%%%%%%%%%%
%
% Stuff for getting the name/document date/title across the header
\makeatletter
\RequirePackage{fancyhdr}
\pagestyle{fancy}
\fancyfoot[C]{\ifnum \value{page} > 1\relax\thepage\fi}
\fancyhead[L]{\ifx\@doclabel\@empty\else\@doclabel\fi}
\fancyhead[C]{\ifx\@docdate\@empty\else\@docdate\fi}
\fancyhead[R]{\ifx\@docauthor\@empty\else\@docauthor\fi}
\headheight 15pt

\def\doclabel#1{\gdef\@doclabel{#1}}
\doclabel{Use {\tt\textbackslash doclabel\{MY LABEL\}}.}
\def\docdate#1{\gdef\@docdate{#1}}
\docdate{Use {\tt\textbackslash docdate\{MY DATE\}}.}
\def\docauthor#1{\gdef\@docauthor{#1}}
\docauthor{Use {\tt\textbackslash docauthor\{MY NAME\}}.}
\makeatother

% Shortcuts for blackboard bold number sets (reals, integers, etc.)
\newcommand{\Reals}{\ensuremath{\mathbb R}}
\newcommand{\Nats}{\ensuremath{\mathbb N}}
\newcommand{\Ints}{\ensuremath{\mathbb Z}}
\newcommand{\Rats}{\ensuremath{\mathbb Q}}
\newcommand{\Cplx}{\ensuremath{\mathbb C}}
%% Some equivalents that some people may prefer.
\let\RR\Reals
\let\NN\Nats
\let\II\Ints
\let\CC\Cplx

%\textbf{Code:}
%\begin{center}
%  \lstinputlisting{NewtonsMethodP5.m}
%\end{center}
%
%\textbf{Console:}
%\begin{center}
%  \lstinputlisting{P5C.txt}
%\end{center}
%\vspace{.15in}


%\begin{figure}[H]
%  \begin{center}
%    \caption{The one-norm unit ball}
%    \includegraphics[width=.76\textwidth]{1norm.png}
%  \end{center}
%\end{figure}




%%%%%%%%%%%%%%%%%%%%%%%%%%%%%%%%%%%%%%%%%%%%%%%%%%%%%%%%%%%%%%%%%%%%%%%%%%%%%%%%%%%%%%%
%%%%%%%%%%%%%%%%%%%%%%%%%%%%%%%%%%%%%%%%%%%%%%%%%%%%%%%%%%%%%%%%%%%%%%%%%%%%%%%%%%%%%%%
% 
% The main document start here.

% The following commands set up the material that appears in the header.
\doclabel{Math 661: Homework 8}
\docauthor{Stefano Fochesatto}
\docdate{\today}


\begin{document}

\begin{exercise}{4.4} Let $M$ be a positive-definite matrix and let 
  \begin{equation*}
    p = -M^{-1}\nabla f(x_k).
  \end{equation*}
  Prove that $p$ is a descent direction for $f$ at $x_k$. 
  \begin{proof} Suppose $p = -M^{-1}\nabla f(x_k)$. For $p$ to be a descent direction we need $p^T \nabla f(x_k) < 0$. By
    substitution we get $(-M^{-1}\nabla f(x_k))^T\nabla f(x_k) = -\nabla f(x_k)(M^{-1})^T\nabla f(x_k)$. Since $M$ is positive definite 
    we know that $\nabla f(x_k)(M^{-1})^T\nabla f(x_k) > 0$ and therefore $-\nabla f(x_k)(M^{-1})^T\nabla f(x_k) < 0$. So $p$ must be a descent direction. 
  \end{proof}
\end{exercise}
\vspace{.5in}

\begin{exercise}{5.2} Let  
  \begin{equation*}
    f(x_1, x_2) = 2x_1^2 + x_2^2 - 2x_1x_2 + 2x_1^3 + x_1^4
  \end{equation*}
  \begin{enumerate}
    \item[(i)]Suppose that the function is minimized stating from $x_0 = (0, -2)^T$. Verify that 
    $p_0 = (0, 1)^T$ is a direction of descent. 
    \begin{proof} First we need to compute the gradient, 
      \begin{equation*}
        \nabla f(x_1, x_2) = (4x_1 - 2x_2 + 6x_1^2 + 4x_1^3, 2x_2 - 2x_1)^T
      \end{equation*}
      Checking if $p_0$ is a descent direction, 
      \begin{equation*}
        p_0^T\nabla f(0, -2) = (0, 1)(4(0) - 2(-2) + 6(0)^2 + 4(0)^3, 2(-2) - 2(0))^T = (0, 1)(4, -4)^T = -4 < 0.
      \end{equation*}
      So $p_0$ is a descent direction. 
    \end{proof}
\vspace{.25in}


\item[(ii)] Suppose that a line search is used to minimize the function $F(\alpha) = f(x_0 + \alpha p_0)$, 
and that a backtracking line search is used to find the optimal step length $\alpha$. Does $\alpha = 1$ 
satisfy the sufficient decrease condition for $\mu = .5$? For what values of $\mu$ does $\alpha = 1$ satisfy the 
sufficient decrease condition?
\begin{proof} Recall that the sufficient decrease condition states that, 
  \begin{equation*}
    f(x_k + \alpha_k p_k) \leq f(x_k) + \mu \alpha p_k^T\nabla f(x_k), 
  \end{equation*}
  for some scalar $0 < \mu< 1$. 
  Checking if $\alpha = 1$ satisfies the sufficient decrease condition for $\mu = .5$, 
  \begin{equation*}
    f(x_0 + (1)p_0) = f(0, -1) = 2(0)^2 + (-1)^2 - 2(0)(-1) + 2(0)^3 + (0)^4 = 1, 
  \end{equation*}
  \begin{equation*}
    f(x_0) + (.5) (1) p_0^T\nabla f(x_0) = 4 + (.5)(1)(-4) = 2.
  \end{equation*}
  So $\alpha = 1$ does not satisfy the sufficient decrease condition for $\mu = .5$. 

  Setting $\alpha = 1$, and solving for $\mu$, 
  \begin{align*}
    f(x_0 + \alpha p_0) \leq f(x_0) + \mu \alpha p_0^T\nabla f(x_0)\\
    1 \leq 4 + \mu(-4)\\
    -3 \leq -4\mu\\
    \frac{3}{4} \geq \mu
  \end{align*}
  So when $\alpha = 1$ in order for $\mu$ to satisfy the sufficient decrease condition we need $\mu \leq \frac{3}{4}$ . 
\end{proof}
  \end{enumerate}
\end{exercise}
\vspace{.5in}

\begin{exercise}{5.5} Let $f$ be the differentiable function that is being minimized with an optimization 
  algorithm that uses an exact line search. Prove that the current search direction is orthogonal to the gradient 
  at the new point $x_{k + 1}$.
  \begin{proof} To show that the current search direction $p_k$ is orthogonal to the gradient at the new point we need to show that 
    $p_k^T\nabla f(x_{k + 1})$. By definition exact line search corresponds to choosing an $\alpha > 0$ which minimizes the function 
    $F(\alpha) = f(x_k + \alpha p_k)$. For $F(\alpha)$ to be a minimizer, the first order necessary condition for optimality requires that 
    $F'(\alpha) = 0$. We showed in class (by a definition of derivative and Taylor Series argument) that $F'(\alpha) = p_k^T\nabla f(x_{k} + \alpha p_k) = p_k^T\nabla f(x_{k + 1}) = 0$.
  \end{proof}
  
\end{exercise}
\vspace{.5in}




\begin{exercise}{2.3} Consider the problem, 
  \begin{equation*}
    \text{ minimize } f(x_1, x_2) = x_1^2 + 2x_2^2.
  \end{equation*}
  \begin{enumerate}
    \item[(i)] If the starting point is $x_0 = (2, 1)^T$, show that the sequence of points generated by the steepest-descent algorithm is given by, 
    \begin{equation*}
      x_k = \frac{1}{3}^k \begin{pmatrix}
        2\\
        (-1)^k
      \end{pmatrix}
    \end{equation*}
    if an exact line search is used. 
    \begin{proof}First let's compute the gradient, 
      \begin{equation*}
        \nabla f(x) = (2x_1,  4x_2).
      \end{equation*}
      As we showed in the previous problem exact line search produces an $\alpha$ such that $p_k^T\nabla f(x_k + \alpha p_k) = 0$. 
      Since we are using steepest descent, we plug in $p_k = -\nabla f(x_k)$, and solving for $\alpha$ find that, 
      \begin{align*}
        p_k^T\nabla f(x_k + \alpha p_k) &= 0\\
        (-\nabla f(x_k))^T\nabla f(x_k + \alpha p_k) &= 0\\
        (-2x_1,  -4x_2)^T(2(x_1 - \alpha(2x_1)), 4(x_2 - \alpha(4x_2))) &= 0\\
        (-2x_1)2(x_1 - \alpha(2x_1)) + (-4x_2)4(x_2 - \alpha(4x_2)) &= 0\\
        (-4x_1)(x_1 - \alpha(2x_1)) + (-16x_2)(x_2 - \alpha(4x_2)) &= 0\\
        -4x_1^2 + \alpha(8x_1^2) - 16x_2^2 + \alpha(64x_2^2) &= 0\\
        \alpha(8x_1^2) + \alpha(64x_2^2) - 4x_1^2 - 16x_2^2  &= 0\\
        \alpha(8x_1^2 + 64x_2^2) - 4x_1^2 - 16x_2^2  &= 0\\
        \alpha  &= \frac{8x_1^2 + 64x_2^2}{4x_1^2 +16x_2^2}\\
        \alpha  &= \frac{x_1^2 + 4x_2^2}{2(x_1^2 + 8x_2^2)}
      \end{align*}
      Since our starting point is $x_0 = (2, 1)^T$ we get an $\alpha_0$ of, 
      \begin{equation*}
        \alpha_0 = \frac{(2)^2 + 4(1)^2}{2((2)^2 + 8(1)^2)} = \frac{1}{3}. 
      \end{equation*}
      Since steepest descent is our update we get, 
      \begin{equation*}
        x_1 = \begin{pmatrix}
          1\\
          2
        \end{pmatrix} + \frac{1}{3}\begin{pmatrix}-2(2)\\-4(1)\end{pmatrix} = \frac{1}{3} \left(\begin{pmatrix}6\\ 3\end{pmatrix}  - \begin{pmatrix} 2^2\\ 4\end{pmatrix}\right) = \frac{1}{3} \begin{pmatrix} 2\\ -1\end{pmatrix}.
      \end{equation*}
      For the next step the $\frac{1}{3}$ will cancel out when we update $\alpha$ and since we are raising every term to
      an even power the $-1$ will also be canceled out so we get $\alpha_1 = \frac{1}{3}$. Computing the next update we get, 
      \begin{equation*}
        x_2 = \frac{1}{3}\begin{pmatrix} 2\\ -1\end{pmatrix} + \frac{1}{3}\left(\frac{1}{3} \begin{pmatrix}-2(2)\\ -4(-1)\end{pmatrix}\right) = \frac{1}{3} \left(  \begin{pmatrix} 2\\ -1\end{pmatrix} + \frac{1}{3} \begin{pmatrix}-2(2)\\ -4(-1)\end{pmatrix}\right) = \frac{1}{3}\left( \frac{1}{3} \begin{pmatrix} 2\\ 1\end{pmatrix} \right) = \left(\frac{1}{3}\right)^2 \begin{pmatrix}2\\ (-1)^2\end{pmatrix}
      \end{equation*}
      We will proceed by induction, suppose that for some $k$ 
      \begin{equation*}
      x_{k} = \frac{1}{3}^k \begin{pmatrix}
        2\\
        (-1)^k
      \end{pmatrix}
    \end{equation*}
    Since steepest descent is our update we know that $x_{k + 1} = x_k + \alpha_k p_k$ where $\alpha_k$ is exact line search and $p_k = -\nabla f(x_k)$. 
    Applying our induction hypothesis and our previous calculation we know that $\alpha_k = \frac{1}{3}$ for the same reasons as in the base case. Thus we get the following,
    \begin{equation*}
      x_{k + 1} = \frac{1}{3}^k \begin{pmatrix}
        2\\
        (-1)^k
      \end{pmatrix}
       + 
       \frac{1}{3}
       \left(
       \frac{1}{3}^k \begin{pmatrix}
        -2(2)\\
        -4((-1))^k
      \end{pmatrix}
      \right)=
      \frac{1}{3}^{k+1} 
      \left(\begin{pmatrix}
        6\\
        3(-1)^k
      \end{pmatrix}
       + 
       \begin{pmatrix}
        -2(2)\\
        -4((-1))^k
      \end{pmatrix}
      \right) = 
      \frac{1}{3}^{k + 1} \begin{pmatrix}
        2\\
        (-1)^{k + 1}
      \end{pmatrix}
    \end{equation*}
    \end{proof}
    \vspace{.15in}

    \item[(ii)] Show that $f(x_{k + 1}) = f(x_k)/9$.
    \begin{proof} By the previous result we know that 
      \begin{equation*}
        x_{k + 1} = 
        \frac{1}{3}^{k + 1} \begin{pmatrix}
          2\\
          (-1)^{k + 1}
        \end{pmatrix}
      \end{equation*}
      Plugging into our equation we get, 
      \begin{equation*}
        f(x_{k + 1}) = \left(\frac{2}{3^{k + 1}}\right)^2 + 2\left(\frac{(-1)^{k + 1}}{3^{k + 1}}\right)^2
      \end{equation*}
    \end{proof}
    We also know that, 
    \begin{equation*}
      f(x_{k}) = \left(\frac{2}{3^{k}}\right)^2 + 2\left(\frac{(-1)^{k}}{3^{k}}\right)^2.
    \end{equation*}
    Through some algebra we find that, 
    \begin{align*}
      f(x_{k + 1}) &= \left(\frac{2}{3^{k + 1}}\right)^2 + 2\left(\frac{(-1)^{k + 1}}{3^{k + 1}}\right)^2,\\
      &= \frac{2^2}{3^{2k + 2}} + 2\left(\frac{(-1)^{2k + 2}}{3^{2k + 2}}\right),\\
      &= \frac{1}{9}\left(\frac{2^2}{3^{2k}} + 2\left(\frac{(-1)^{2k}}{3^{2k}}\right)\right),\\
      &= \frac{1}{9}\left(\left(\frac{2}{3^{k}}\right)^2 + 2\left(\frac{(-1)^{k}}{3^{k}}\right)^2\right),\\
      &= \frac{1}{9}\left(f(x_k)\right).
    \end{align*}
    \vspace{.15in}


    \item[(iii)] Compare the results in $(ii)$ to the bounds on the convergence rate of the steepest descent method 
    when minimizing a quadratic function. What conclusions can you draw regarding this method?
    \begin{proof}[Solution] Recall from Section 12 that the bounds on the convergence of the steepest descent method 
      when minimizing a quadratic function using exact line search, is given by the following, 
      \begin{align*}
        f(x_{k + 1}) - f(x_*) &\leq \left[ \frac{cond(Q) - 1}{cond(Q) + 1}\right]^2(f(x_{k}) - f(x_*))\\
        \frac{f(x_{k + 1}) - f(x_*)}{f(x_{k}) - f(x_*)} &\leq \left[ \frac{cond(Q) - 1}{cond(Q) + 1}\right]^2
      \end{align*}
      For our problem we had a, 
      \begin{equation*}
        Q = \begin{pmatrix}
          2 & 0\\
          0 & 4
        \end{pmatrix}
      \end{equation*}
      Computing the convergence constant we get, 
      \begin{equation*}
        \frac{f(x_{k + 1}) - f(x_*)}{f(x_{k}) - f(x_*)} &\leq \left[ \frac{cond(Q) - 1}{cond(Q) + 1}\right]^2 = \left(\frac{2 - 1}{2 + 1}\right)^2 = \frac{1}{9}. 
      \end{equation*}
      Our result is exactly what we expect from part $(ii)$, confirming again that this method converges linearly, and for this problem we 
      have a rate constant of $1/9$. 


    \end{proof}

  \end{enumerate}
  
  
\end{exercise}
\vspace{.5in}



\begin{exercise}{3.7} Let $B_{k + 1}$ be obtained from $B_k$ using the update formula, 
  \begin{equation*}
    B_{k + 1} = B_k + \frac{(y_k - b_ks_k)v^T}{v^Ts_k}
  \end{equation*}
  where $v$ is a vector such that $v^Ts_k \neq 0$. Prove that $B_{k + 1}s_k = y_k$. 
  \begin{proof} Consider $B_{k + 1}s_k$ and by substitution we get, 
    \begin{align*}
      B_{k + 1}s_k &= \left(B_k + \frac{(y_k - B_ks_k)v^T}{v^Ts_k}\right)s_k,\\
     &= B_k s_k + \frac{(y_k - B_ks_k)v^Ts_k}{v^Ts_k},\\
     &= B_k s_k + \frac{(y_k - B_ks_k)(v^Ts_k)}{v^Ts_k},\\
     &= B_k s_k + (y_k - B_ks_k),\\
     &= y_k.
    \end{align*}
    
  \end{proof}
\end{exercise}
\vspace{.5in}


\begin{exercise}{P16} Suppose $x_k \in \RR^n$ is any iterate, and suppose that $\nabla f(x_k)$ is a nonzero vector. Compute all $p\in \RR^n$
  which solve the problem, 
  \begin{equation*}
    \min_{p \neq 0} \frac{p^T \nabla f(x_k)}{||p|| ||\nabla f(x_k)||}
  \end{equation*}
  \begin{proof}[Solution] Note that $\nabla f(x_k)$ and $p$ are nonzero vectors and since we know that the $p^T \nabla f(x_k) = ||p||||\nabla f(x_k)|| \cos(\theta)$
    where $\theta$ is the angle between vectors $p$ and $\nabla f(x_k)$ by substitution our objective function turns into, 
    \begin{equation*}
      \min_{p \neq 0} \frac{||p||||\nabla f(x_k)|| \cos(\theta)}{||p|| ||\nabla f(x_k)||} = \cos(\theta)
    \end{equation*}
    Therefore $p$ must be a vector orthogonal to $\nabla f(x_k)$. Kind of a roundabout way of saying every $p\in \RR^n$ lies on the hyper plane defined by $p^T \nabla f(x_k) = 0$. 
  \end{proof}
\end{exercise}
\vspace{.5in}

\begin{exercise}{P17} Consider the one-variable problem
  \begin{equation*}
    min_{x \in \RR}f(x)
  \end{equation*}
  where $f: \RR \to \RR$ is twice continuously-differentiable. Recall that Newton method for the above minimization solves $f'(x) = 0$
  by the formulas $p_k =-f'(x_k)/f''(x_k)$ and $x_{k+1} = x_k + p_k$. 
  
  The secant method for minimization only differs from the Newton method by replacing the second derivative with a difference quotient approximation based on 
  the last two iterates:
  \begin{equation*}
    f''(x_k) \approx \frac{f'(x_k) - f'(x_{k - 1})}{x_k - x_{k - 1}}.
  \end{equation*}
  Thus the secant method computes the step search vector by, 
  \begin{equation*}
    p_k = -\frac{(x_k - x_{k - 1}) f'(x_k)}{f'(x_k) - f'(x_{k - 1})}
  \end{equation*}
  with $x_{k + 1}= x_k + p_k$ as before. 
  \begin{enumerate}
    \item[(a)] Implement the secant method. Include a stopping criterion $|f'(x_k)|< tol$.
    \begin{proof}[Solution]
      \textbf{Code:}
      \begin{center}
        \lstinputlisting[basicstyle = \footnotesize]{secant}
      \end{center}
    \end{proof}
\vspace{.5in}


    \item[(b)] Use your code to accurately solve $min_{x \in \RR}f(x)$, e.g. with $tol = 10^{-8}$, 
    for the following functions and initial iterates:
    \begin{enumerate}
      \item[(i)] $f(x) = x^3 - 2sin(x)$ with $x_0 = 0$ and $x_1 = 1$
      \begin{proof}[Solution]

        \textbf{Console:}
        \begin{center}
          \lstinputlisting[basicstyle = \footnotesize]{r1.txt}
        \end{center}
      \end{proof}
      \vspace{.15in}

    

      \item[(ii)] $f(x) = 3x^4 - 4x^3 + 3x^2 - 6x$ with $x_0 = -1$ and $x_1 = 0$.   
      \begin{proof}[Solution] 

        \textbf{Console:}
        \begin{center}
          \lstinputlisting[basicstyle = \footnotesize]{r2.txt}
        \end{center}
      \end{proof}
    \end{enumerate} 
\vspace{.5in}

\item[(c)] In part $(ii)$ above the exact minimum is at $x_* = 1$. Compute the errors $e_k = x_k - x_*$. Give evidence that the convergence is 
superlinear. Using the notation of section 2.5, what is you estimate of the rate exponent r?.
\begin{proof}[Solution] Computing $e_k$ we get the following, \\

  \textbf{Console:}
        \begin{center}
          \lstinputlisting[basicstyle = \footnotesize]{r3.txt}
        \end{center}

Setting $r = 1$ and computing $\frac{||e_{k + 1}||}{||e_k||}$ we can see the sequence is heading towards zero, which by definition is evidence of 
superlinear convergence.\\

\textbf{Console:}
\begin{center}
  \lstinputlisting[basicstyle = \footnotesize]{r4.txt}
\end{center}


I don't have a good sense of how to estimate $r$ from just looking at our list of errors, however I found a paper\cite{1} which gives a formula 
for estimating r, 
\begin{equation*}
  r \approx \dfrac{log|(x_{n + 1} - x_n)/(x_{n} - x_{n - 1})|}{log|(x_{n} - x_{n - 1})/(x_{n - 1} - x_{n - 2})|}
\end{equation*}
Computing this sequence we get the following, \\

\textbf{Console:}
\begin{center}
  \lstinputlisting[basicstyle = \footnotesize]{r5.txt}
\end{center}

Towards the end it seems to be converging to $\sim$ 1.6. Double checking by computing $\frac{||e_{k + 1}||}{||e_k||^{1.6}}$, if 
we have a good estimate for $r$ we should expect this sequence to converge to some $0< C < \infty$.\\

\textbf{Console:}
\begin{center}
  \lstinputlisting[basicstyle = \footnotesize]{r6.txt}
\end{center}
Since this sequence seems to be staying at around $1$ it seems that $r = 1.6$ is a good estimate for $r$ especially since we 
discussed it in class and you told us secant method converges with $r$ as the golden ratio. 

\end{proof}
  \end{enumerate}
  
\end{exercise}
\begin{thebibliography}{1}  % "2" because there are two references
  \bibitem{1}
  Senning, J.R., n.d. COMPUTING AND ESTIMATING THE RATE OF CONVERGENCE.
  \end{thebibliography}
\end{document}


https://www.maa.org/press/periodicals/loci/joma/iterative-methods-for-solving-iaxi-ibi-gauss-seidel-method
    https://www.maa.org/press/periodicals/loci/joma/iterative-methods-for-solving-iaxi-ibi-jacobis-method\\




\end{document}




























